\section{Discussion}
\label{sec:discussion}

\subsection{Physical interpretation of performance gains}

The 3,500-case paired comparison directly connects theoretical failure mechanisms to concrete quantitative improvements. MIKU raises collision-free goal reaching from 62.69\% to 77.46\% against baseline B0, yielding a paired increase of 14.77~pp with a 95\% confidence interval of 13.54 to 15.97. This performance gain is concentrated in configurations governed by severe spatial or temporal bottlenecks:
under narrow multi-obstacle geometry, where greedy obstacle-wise decisions cross each other and collapse opposing path bounds such that $l^+ < l^-$, MIKU's group-level partition restores a continuous, feasible center corridor, increasing success by 65.4 pp;
under delayed-crossing geometry, where time-blind bounds falsely treat future arrival-time occupancies as immediate blockages, arrival projection and temporal homotopy preserve reachable passing windows before and after the crossing obstacle, increasing success by 26.8 pp.
In contrast, in scenarios naturally characterized by a single dominant corridor, such as simple pedestrian crossings and high-speed vehicle cut-ins, performance remains consistent with the baseline, verifying that MIKU specifically targets and resolves diagnosed interface bottlenecks.

Ablation experiments substantiate this mechanism attribution. Removing Top-$K$ spatial homotopy under ablation~A3 or the threat-aware margin under ablation~A4 reduces success by 10.17~pp and 10.74~pp, respectively, demonstrating that group-level lateral coordination provides the foundational clearance necessary for spatial feasibility. Removing the temporal homotopy graph under ablation~A5 selectively cuts delayed-crossing success from 89.6\% to 66.8\%, a loss of 22.8~pp that confirms its decisive role when multiple temporal passages coexist. Finally, removing the robust occupancy tube under ablation~A6 causes prediction-noise collision rates to jump from 3.6\% to 11.0\%, affirming that uncertainty-aware buffer expansion is essential for absorbing real-world tracking and perception noise \citep{yoon2024stcorridor,yangbo2023ped,pathspeedstructured2024}.

These findings align with prior decomposition-based planning literature while clarifying the locus of improvement. \citet{fan2018baidu} and \citet{zhang2020pathqp} demonstrated that piecewise-jerk QP formulations achieve production-grade smoothness and latency. Corridor-based approaches \citep{liu2017sfc,yoon2024stcorridor} showed that explicit feasible-region construction improves trajectory quality in structured environments; MIKU extends this principle by coupling spatial and temporal corridor selection through a shared homotopy graph. The iterative anchoring strategy of \citet{chen2019em} feeds timing information back between stages, yet our B2 results, reaching 62.09\% success and remaining comparable to the time-blind B0, confirm that iterative arrival-time updates alone cannot recover lateral classes discarded by an irreversible initial side decision. For traffic system designers, the practical consequence is that upgrading the constraint interface layer---without replacing the downstream QP solvers---can recover a substantial fraction of the performance gap between decomposed and joint planners while preserving the computational budget required for real-time deployment.

\subsection{Computational efficiency and engineering trade-offs}

In operational autonomous driving systems, planning algorithms must operate within strict onboard latency budgets, typically 50 to 100~ms. Across 3,500 evaluated instances, MIKU achieves a median runtime of 9.98 ms and a P95 of 55.93 ms, demonstrating that interaction-aware constraint reconstruction adds negligible overhead during nominal driving conditions.

Computational cost concentrates primarily in the distribution tail: complete-call P99 latency rises to 221.80~ms compared with 138.88~ms for B0. This tail behavior reflects the physical reality of dense multi-obstacle traffic: branch expansion in the spatial partition and temporal homotopy graph only activates when multiple obstacles simultaneously compete for shared roadway space. For production deployment, timing guards should be budgeted against the P99 tail rather than median latency. Furthermore, under prediction noise, the 7.0 pp collision reduction incurs a modest 1.12 pp progress reduction and a 0.041 m/s\textsuperscript{3} jerk increase, illustrating an intentional and controllable engineering trade-off between absolute safety margins and aggressive forward progress.

\subsection{Comparison with joint search and iterative baselines}

Under the 70-case finite-grid protocol, the joint spatiotemporal reference B3 achieves 80.0\% success compared with 77.1\% for MIKU, with a paired difference of $-2.86$~pp, a 95\% confidence interval of $[-11.43,5.71]$, and $p=0.754$. B3's median runtime is 1,835.43 ms compared with MIKU's 10.44 ms, so the finite-grid search trades a modest success increment for a much higher computational cost.

Conversely, the iterative PVD baseline B2 executes three sequential path--speed iterations, requiring a median runtime of 31.20~ms, yet achieves only 62.09\% success, which is comparable to B0's 62.69\% and 15.37~pp lower than MIKU.

\subsection{Comprehensive insights from multi-protocol evaluation}

The multi-tiered experimental framework evaluates planning performance across distinct operational dimensions. The open-loop 3,500-scenario protocol provides statistical power and mechanism isolation under controlled environmental distributions. The 700-scenario closed-loop rolling replanning protocol demonstrates that open-loop advantages reliably accumulate over extended operational horizons, lifting overall goal-reaching success from 62.29\% to 72.29\%, an improvement of 10.0~pp, while maintaining trajectory stability across successive planning cycles.

Finally, integration into Apollo 11.0 validates architectural compatibility on an established production platform. Across dynamic overtaking, narrow passage navigation, and lane diversion scenarios in Dreamview, MIKU successfully drives the vehicle without modifying existing quadratic programming solvers or control interfaces.
